\documentclass[12pt]{article}
\usepackage[T1]{fontenc}
\usepackage[utf8]{inputenc}
\usepackage{lmodern}
\usepackage{textcomp}
\usepackage{enumitem}
\usepackage{geometry}
\usepackage{graphicx}
\usepackage{amsmath, amssymb}
\geometry{margin=1in}
\begin{document}
\begin{center}
Biology - Undergraduate Level Assessment 22 \
Total Marks: 40 \
Answer all questions.
\end{center}

\section*{Multiple Choice (10 marks)}
\begin{enumerate}[label=\arabic*.]
\item Which of the following adaptive features would least likely be found in an animal living in a hot arid environment?
\begin{enumerate}[label=(\alph*)]
    \item High metabolic rate to generate heat
    \item Storage of water in fatty tissues
    \item Long loops of Henle to maximize water reabsorption
    \item Short loops of Henle to maximize water secretion
    \item Large amounts of body hair to trap heat
\end{enumerate}
\item During the courtship of the common tern, the male presents afish to another tern. Account for this behavior and discuss whycourtship in general may benecesaryfor reproduction.
\begin{enumerate}[label=(\alph*)]
    \item Courtship is primarily for entertainment purposes among terns
    \item Courtship is not necessary for reproduction
    \item Courtship is a way to show dominance
    \item Courtship is a random behavior with no impact on reproduction
    \item Courtship is necessary for physical transformation required for reproduction
\end{enumerate}
\item What function does oxygen serve in the body?
\begin{enumerate}[label=(\alph*)]
    \item Oxygen is used to maintain the temperature of cells.
    \item Oxygen is used by cells to produce alcohol.
    \item Oxygen is used to produce carbon dioxide in cells.
    \item Oxygen is solely responsible for stimulating brain activity.
    \item Oxygen is used to synthesize proteins in cells.
\end{enumerate}
\item Which of the following is the source of oxygen produced during photosynthesis?
\begin{enumerate}[label=(\alph*)]
    \item O 2
    \item CH 4 (Methane)
    \item CO 2
    \item H 2 O
    \item N 2 (Nitrogen)
\end{enumerate}
\item What is the genetic basis of Down's syndrome?
\begin{enumerate}[label=(\alph*)]
    \item Duplication of chromosome 14
    \item Translocation between chromosomes 14 and 21
    \item Deletion in chromosome 21
    \item Mutation in chromosome 14
    \item Duplication of chromosome 18
\end{enumerate}
\end{enumerate}

\section*{True / False (10 marks)}
\textit{Answer True or False}
\begin{enumerate}[label=\arabic*.]
\setcounter{enumi}{5}
\item True or False: The light-dependent reactions produce NADP+, ADP, and sugar during photosynthesis.
\item True or False: Comparing gene frequencies and genotype frequencies over two generations can determine if a population is in genetic equilibrium.
\item True or False: The probability of getting exactly three heads in five flips of a balanced coin is 12.5 percent.
\item True or False: A photosynthetic plant in nutrient-deficient soil typically exhibits dark green, robust growth with normal height and vigor.
\item True or False: Cholesterol uptake through low density lipoprotein complexes does not involve adhesion plaques.
\end{enumerate}

\section*{Long Form Response (20 marks)}
\textit{Answer each question with a clear and complete explanation.}
\begin{enumerate}[label=\arabic*.]
\setcounter{enumi}{10}
\item Discuss the genetic inheritance of hemophilia in a phenotypically normal couple who have a son with the disorder. What is the probability that their next child, if a girl, will also have hemophilia?
\item Discuss the concept of coevolution, particularly in the context of predator-prey relationships, and analyze its implications for biodiversity and ecosystem dynamics.
\end{enumerate}

\vfill
\noindent\textit{End of Paper}
\end{document}